% !PRINTED none
% Abschnitt (part) title page: "Vierzehnter Abschnitt." / section title, rule, then opening of §.152 with its heading and the start of the section's running text and first footnote. No running-head page number is printed on this title page.
\begin{center}
\gesperrt{Fourteenth Section.}
\end{center}

\begin{center}
\gesperrt{Application of the Permutation Groups to Equations.}
\end{center}

\begin{center}
\rule{2cm}{0.4pt}
\end{center}

\begin{center}
\S.\,152.
\end{center}

\begin{center}
\gesperrt{Effect of the Permutation Groups on Functions of Independent Variables.}
\end{center}

For a deeper penetration into algebra, a more exact study of the permutation groups is, after the results of the preceding section, necessary. We first set up some general theorems about them, which form the foundation for what follows\footnote{For the theory of permutations the following are especially to be noted: \gesperrt{Cauchy}, ``Journal de l'École polytechn. X. cah. (1815)'' (several memoirs). ``C. \gesperrt{Jordan}, Traité des substitutions et des équations algébriques'', Paris 1870. \gesperrt{Netto}, ``Substitutionentheorie und ihre Anwendung auf die Algebra.'' Leipzig 1882.}.

Let
\weq{1}{u,\ u_1,\ u_2 \ldots u_{m-1}}
be a system of mutually independent symbols (variables), and
\weq{2}{\psi(u,\ u_1,\ u_2 \ldots u_{m-1})}
an integral rational function of them with coefficients from an arbitrary field.

Two such functions $\psi$ are to be regarded as equal only when, in the expressions ordered according to powers and products of the $u$, corresponding terms have the same coefficients.
